\subsection{Crandall's complete accelerated programme}
\label{subsec:Crandall-complete-audit}

This subsection reconstructs all twelve printed pages of Crandall's 1978
article from the published witness \cite[pp.~1281--1292]{Crandall1978}.  The
source uses only the accelerated odd map \eqref{eq:crandall-C}.  Its
trajectory is the indexed sequence
\[
 \mathcal T_m=\bigl(C_{\mathrm{Cr}}^j(m)\bigr)_{j\geq1},
\]
stopped after the least index $j$ for which $C_{\mathrm{Cr}}^j(m)=1$, if
such an index exists, and left infinite otherwise.  Its height is that least
index, or $\infty$ when there is none.  Thus the braces used in the source's
display denote an ordered sequence, not its underlying set; repeated values
in a nonterminating periodic trajectory do not make its height finite.  In
particular, the source table has $\mathcal T_1=(1)$ and $h(1)=1$; this is not
the zero-step hitting-time convention.

For a segment $x_i=C_{\mathrm{Cr}}^i(x_0)$, put
\begin{equation}
\label{eq:Crandall-cumulative-time}
 e_i=\vtwo(3x_i+1),
 \qquad
 A_0=0,
 \qquad
 A_j=\sum_{i=0}^{j-1}e_i.
\end{equation}
Repeated use of Proposition~\ref{prop:first-return-clocks} gives
\begin{equation}
\label{eq:Crandall-three-clocks}
 C_{\mathrm{Cr}}^j(x_0)=T^{A_j}(x_0)=U^{A_j+j}(x_0).
\end{equation}
Thus every height, trajectory position, and period in Crandall's paper is
accelerated time $j$.  Neither $A_j$ nor $A_j+j$ may be substituted for it.

\subsubsection{Large finite excursions and the random-walk nonclaim}

\begin{proposition}[Crandall, Theorem 2.2]
\label{prop:Crandall-excursions}
For $k\geq2$, let $m_k=2^k-1$.  Then
\begin{equation}
\label{eq:Crandall-excursion-orbit}
 C_{\mathrm{Cr}}^j(m_k)=3^j2^{k-j}-1
 \qquad(0\leq j<k),
\end{equation}
and hence
\begin{equation}
\label{eq:Crandall-excursion-ratio}
 \frac{\sup\mathcal T_{m_k}}{m_k}>
 \left(\frac32\right)^{k-1}.
\end{equation}
Consequently the family of ratios
\[
 \left\{\frac{\sup\mathcal T_m}{m}:m\in\Opos\right\}
\]
is unbounded.  For the ratios with denominator $m^\alpha$, this construction
proves unboundedness only for $\alpha<\log_2 3$.
\end{proposition}

\begin{proof}
Equation~\eqref{eq:Crandall-excursion-orbit} is immediate at $j=0$.  If it
holds with $j<k-1$, then
\[
 3(3^j2^{k-j}-1)+1
 =2(3^{j+1}2^{k-j-1}-1),
\]
whose second factor is odd; the next accelerated exponent is therefore one.
At $j=k-1$, the trajectory contains $2\cdot3^{k-1}-1$.  Direct
cross-multiplication gives
\[
 \frac{2\cdot3^{k-1}-1}{2^k-1}>
 \frac{3^{k-1}}{2^{k-1}}
\]
because $3^{k-1}>2^{k-1}$.  This proves the stated bound.  Dividing the same
member by $m_k^\alpha$ yields an exponential factor
$3^{k-1}/2^{\alpha k}$, which this construction forces upward exactly when
$\alpha<\log_2 3$ \cite[pp.~1282--1283, Thm.~2.2 and following paragraph]
{Crandall1978}.
\end{proof}

Section~3 assigns the heuristic probabilities
$\Pr(e(m)=a)=2^{-a}$ for $a\geq1$.  Its model logarithmic drift is
\begin{equation}
\label{eq:Crandall-heuristic-drift}
 \sum_{a\geq1}2^{-a}\log(3/2^a)=-\log(4/3).
\end{equation}
It then proposes the individual estimate
$h(m)\sim\log m/\log(4/3)$.  For the source average
\begin{equation}
\label{eq:Crandall-average-height}
 H(x)=\frac2x\sum_{\substack{m\leq x\\m\in\Opos}}h(m),
\end{equation}
Conjecture~3.1 is
\begin{equation}
\label{eq:Crandall-average-conjecture}
 H(x)\sim\frac{2\log x}{\log(16/9)}.
\end{equation}

\begin{nonclaim}
\label{nonclaim:Crandall-random-walk}
Equations \eqref{eq:Crandall-heuristic-drift}--
\eqref{eq:Crandall-average-conjecture} are explicitly heuristic or
conjectural.  They prove neither termination nor unboundedness.  The source
also reports Everett's almost-everywhere descent theorem, but does not prove
it.  Everett's primary proof and the exact variable-clock transfer are now
reconstructed in Theorem~\ref{thm:Everett-density-one-descent} and
Proposition~\ref{prop:Everett-Crandall-crosswalk}; neither follows from the
heuristic drift displayed here.  The sentence after Conjecture~3.1 calls it
``(3.1)'', duplicating the number of the individual-height heuristic
\cite[pp.~1282--1284]{Crandall1978}.
\end{nonclaim}

\subsubsection{Backward coordinates and the exact finite-height address}

Crandall defines, for $a\in\Z_{>0}$ and $n\in\Q$,
\begin{equation}
\label{eq:Crandall-backward-branch}
 B_a(n)=\frac{2^a n-1}{3},
 \qquad
 B_{a_j\cdots a_1}=B_{a_j}\circ\cdots\circ B_{a_1}.
\end{equation}
The subscript order is backward: when all displayed values are integral,
$C_{\mathrm{Cr}}$ removes the leftmost applied branch $B_{a_j}$.

For a natural-order address $\mathbf a=(a_1,\ldots,a_j)$, define
\begin{equation}
\label{eq:Crandall-address-recursion}
 n_0=1,
 \qquad
 n_i=B_{a_i}(n_{i-1})\quad(1\leq i\leq j).
\end{equation}
Let $G^\sharp$ be the set of such finite addresses satisfying
\begin{align}
 a_1&>2,
 \label{eq:Crandall-address-first-gate}\\
 2^{a_i}n_{i-1}&\equiv1\pmod3
 &&(1\leq i\leq j),
 \label{eq:Crandall-address-integrality-gate}\\
 2^{a_i}n_{i-1}&\not\equiv1\pmod9
 &&(1\leq i<j).
 \label{eq:Crandall-address-continuation-gate}
\end{align}

\begin{theorem}[Crandall's finite-height address theorem, repaired]
\label{thm:Crandall-address-bijection}
Let
\[
 \mathcal M=\{m\in\Opos:m>1,\ h(m)<\infty\}.
\]
The map
\begin{equation}
\label{eq:Crandall-address-evaluation}
 \operatorname{ev}:G^\sharp\longrightarrow\mathcal M,
 \qquad
 (a_1,\ldots,a_j)\longmapsto B_{a_j\cdots a_1}(1)
\end{equation}
is a bijection.  Its inverse sends $m\in\mathcal M$ to the length-$h(m)$
address
\begin{equation}
\label{eq:Crandall-address-inverse}
 a_i=e(C_{\mathrm{Cr}}^{h(m)-i}(m))
 \qquad(1\leq i\leq h(m)).
\end{equation}
The exact trajectory of $n_j=\operatorname{ev}(\mathbf a)$ is
\[
 n_j\longmapsto n_{j-1}\longmapsto\cdots\longmapsto n_1
 \longmapsto1,
\]
so its height is exactly $j$.
\end{theorem}

\begin{proof}
Suppose first that \eqref{eq:Crandall-address-integrality-gate} holds through
index $i$.  Then $n_i=(2^{a_i}n_{i-1}-1)/3$ is integral.  Since the numerator
is odd and the denominator is odd, $n_i$ is odd.  Moreover
\[
 n_i\equiv0\pmod3
 \quad\Longleftrightarrow\quad
 2^{a_i}n_{i-1}\equiv1\pmod9.
\]
Thus \eqref{eq:Crandall-address-continuation-gate} says exactly that every
nonterminal $n_i$ is prime to $3$, which is necessary for another inverse
branch to be integral.  From
$3n_i+1=2^{a_i}n_{i-1}$ and oddness of $n_{i-1}$, one has
$\vtwo(3n_i+1)=a_i$ and hence $C_{\mathrm{Cr}}(n_i)=n_{i-1}$.  Finally,
$a_1>2$ ensures $n_1\ne1$, so the first occurrence of $1$ is terminal and
the height is $j$.

Conversely, a finite trajectory from $m>1$ supplies the valuations in
\eqref{eq:Crandall-address-inverse}.  Reversing its branch equations gives
\eqref{eq:Crandall-address-recursion}.  Every nonterminal orbit value must be
prime to $3$: if $3\mid n_i$, then no equation
$3u+1=2^a n_i$ is possible modulo $3$.  Hence all three gates hold.  The
determinism of $C_{\mathrm{Cr}}$ recovers every $a_i$ uniquely, so evaluation
and \eqref{eq:Crandall-address-inverse} are literal two-sided inverses.  This
is the corrected content of Lemmas~4.1--4.3 and Theorem~4.1
\cite[pp.~1284--1285]{Crandall1978}.
\end{proof}

\begin{sourceissue}[two finite-address endpoint defects]
\label{sourceissue:Crandall-address-endpoints}
The printed definition of $G$ imposes a non-$1\pmod9$ condition on $a_1$
and a terminal mod-$3$ condition on $a_j$.  At length $j=1$, both apply to
the same letter.  It thereby excludes the valid height-one example
\[
 B_6(1)=21,
 \qquad C_{\mathrm{Cr}}(21)=1,
 \qquad 2^6\equiv1\pmod9.
\]
The mod-$9$ gate belongs only to nonterminal indices, as in
\eqref{eq:Crandall-address-continuation-gate}.  Independently, the printed
Lemma~4.3 includes ``an integer of height $j$'' while $a_1>2$ excludes
$B_2(1)=1$; the source table nevertheless has $h(1)=1$.  Restricting the
lemma to the $m>1$ set used by Theorem~4.1 repairs this endpoint.  Likewise,
inside $\mathcal M$ the height-one family $(4^s-1)/3$ has $s\geq2$, not
merely $s\geq1$ \cite[pp.~1282, 1284--1285]{Crandall1978}.
\end{sourceissue}

The address theorem is genuinely a bijection, not a merely suggestive
parameterization: its domain, codomain, evaluation, inverse, and exact clock
are all displayed.  It is also a finite-orbit theorem.  It provides no
infinite address space and no conclusion about nonterminating positive
orbits.

\subsubsection{Fixed-height and total termination counts}

Let
\[
 \pi_h(x)=\#\{m\in\mathcal M:h(m)=h,\ m\leq x\}.
\]
The elementary size estimate is
\begin{equation}
\label{eq:Crandall-backward-size}
 B_{a_j\cdots a_1}(1)
 <\frac{2^{a_1+\cdots+a_j}}{3^j}.
\end{equation}
It follows by induction from \eqref{eq:Crandall-backward-branch}, retaining
the discarded negative translation rather than replacing the branch by an
equality.

\begin{sourceissue}[domain of the address-counting lemma]
\label{sourceissue:Crandall-lemma52-domain}
Printed Lemma~5.2 says that for every real $z>0$, the number of length-$j$
addresses in $G$ with coordinate sum at most $z$ is at least
\begin{equation}
\label{eq:Crandall-address-count}
 \left(2\left\lfloor\frac{z-2}{6j}\right\rfloor\right)^j.
\end{equation}
At $z=1,j=2$, the right side is $4$, while no two positive coordinates have
sum at most $1$.  The proof's blocks of six require a nonnegative number of
blocks.  The statement is correct at the scope used later after requiring
$z\geq2$, or globally after replacing the floor by its nonnegative part.
\end{sourceissue}

\begin{theorem}[Crandall, fixed accelerated height]
\label{thm:Crandall-fixed-height-count}
There are real constants $r,x_0>0$, independent of $h$, such that for every
$h\in\Z_{>0}$,
\begin{equation}
\label{eq:Crandall-fixed-height-hypothesis}
 x>\max\{x_0,2^{h/r}\}
\end{equation}
implies
\begin{equation}
\label{eq:Crandall-fixed-height-bound}
 \pi_h(x)>\frac{(\log_2(x^r))^h}{h!}.
\end{equation}
In particular, infinitely many positive odd integers have each prescribed
positive accelerated height $h$.
\end{theorem}

\begin{proof}
Theorem~\ref{thm:Crandall-address-bijection},
\eqref{eq:Crandall-backward-size}, and the repaired form of
\eqref{eq:Crandall-address-count} give the source's explicit preliminary
bound
\[
 \pi_h(x)\geq
 \left(\frac{\log_2(3^h x/4)}{3h}-2\right)^h
\]
once the displayed base is positive.  Choose $x_0$ so that $x/4>x^{1/2}$
for $x>x_0$.  Choose $0<r<1$ with
\begin{equation}
\label{eq:Crandall-r-choice}
 r\log_2(3/64)+\frac12>3\mathrm e r>0.
\end{equation}
If also $x>2^{h/r}$, then $h<r\log_2x$.  Substituting this strict inequality
in the preliminary bound and using $h!\mathrm e^h>h^h$ yields
\eqref{eq:Crandall-fixed-height-bound}.  This follows the source proof while
making its strict threshold and the repaired Lemma~5.2 domain explicit
\cite[p.~1286, Thm.~5.1]{Crandall1978}.
\end{proof}

Let
\begin{equation}
\label{eq:Crandall-pi-total}
 \pi(x)=\#\{m\in\mathcal M:m\leq x\}
       =\sum_{h\geq1}\pi_h(x).
\end{equation}

\paragraph{The exact handbook route.}
Crandall's remark after Lemma~6.1 says only that ``various formulas''
pertaining to the lemma occur in reference~4; it gives no handbook section,
page, or formula number \cite[p.~1287]{Crandall1978}.  The relevant formulas
can nevertheless be identified exactly.  Abramowitz--Stegun use $P(z)$ for
the standard normal lower tail, but also use $P(\chi^2\mid\nu)$ and
$Q(\chi^2\mid\nu)=1-P(\chi^2\mid\nu)$ for the lower and upper chi-square
probabilities.  Their formulas 26.2.2--26.2.6 give $P(0)=Q(0)=1/2$; formula
26.2.29 relates this $P$ to $\operatorname{erf}$; formula 26.4.11 gives the
large-$\nu$ normal coordinate
\[
 \frac{\chi^2-\nu}{\sqrt{2\nu}};
\]
and formula 26.4.21 gives the exact finite Poisson sum below
\cite[printed pp.~931, 934, 940--941]{AbramowitzStegun1964}.  These are
identified relevant formulas, not locators supplied by Crandall.

\begin{proposition}[exact Poisson--gamma--chi-square coordinate]
\label{prop:AS-Poisson-gamma-morphism}
For $n\in\Nzero$ and $t\geq0$, set
\begin{equation}
\label{eq:AS-Poisson-gamma-morphism}
 F_n(t)=e^{-t}\sum_{j=0}^{n}\frac{t^j}{j!}.
\end{equation}
Then, with $\Gamma(a,t)$ denoting the upper incomplete gamma function,
\begin{equation}
\label{eq:AS-Poisson-gamma-identities}
 F_n(t)=\frac{\Gamma(n+1,t)}{\Gamma(n+1)}
       =Q(2t\mid2n+2).
\end{equation}
For fixed $n$,
\begin{equation}
\label{eq:AS-Poisson-gamma-derivative}
 F_n'(t)=-e^{-t}\frac{t^n}{n!}\qquad(t>0),
\end{equation}
so $F_n$ is a continuous strictly decreasing bijection
$[0,\infty)\to(0,1]$.  Its inverse is the unique upper-tail quantile on this
fixed-$n$ fibre; no elementary inverse formula is asserted.
\end{proposition}

\begin{proof}
Integration by parts gives
\[
 \Gamma(n+1,t)=n!e^{-t}\sum_{j=0}^{n}\frac{t^j}{j!}
\]
by induction from $\Gamma(1,t)=e^{-t}$, proving the first equality in
\eqref{eq:AS-Poisson-gamma-identities}.  In Abramowitz--Stegun 26.4.21 take
$c=n+1$, $\nu=2c$, $m=t$, and $\chi^2=2t$ to obtain the second equality
literally \cite[printed p.~941, Eq.~26.4.21]{AbramowitzStegun1964}.
Differentiating the finite sum makes every term cancel except the displayed
last term in \eqref{eq:AS-Poisson-gamma-derivative}.  Finally $F_n(0)=1$,
while $F_n(t)\to0$ because an exponential dominates the fixed polynomial;
strict monotonicity supplies the stated fibres and inverse.
\end{proof}

\begin{lemma}[completion of Crandall's Poisson lemma]
\label{lem:Crandall-Poisson-half-mass}
Let $k(t)$ be integer-valued and suppose
\begin{equation}
\label{eq:Crandall-Poisson-moving-cutoff}
 \frac{k(t)-t}{\sqrt t}\longrightarrow0.
\end{equation}
Then
\begin{equation}
\label{eq:Crandall-Poisson-half-mass}
 \lim_{t\to\infty}e^{-t}
 \sum_{u=0}^{k(t)}
 \frac{t^u}{u!}=\frac12.
\end{equation}
Consequently the same limit holds after deleting $u=0$ for each of the
following three exact upper cutoffs:
\begin{equation}
\label{eq:Crandall-Poisson-three-cutoffs}
 k(t)=\lfloor t\rfloor-1,
 \qquad
 k(t)=\lfloor t\rfloor,
 \qquad
 k(t)=\lceil t\rceil-1.
\end{equation}
The first is Crandall's printed strict cutoff
$u<\lfloor t\rfloor$; the third is the maximal integer cutoff uniformly
authorized by the strict height hypothesis in Theorem~5.1.
\end{lemma}

\begin{proof}
The handbook route to the source case is now literal.  Put
$n=\lfloor t\rfloor$, $\nu=2n$, and $\chi^2=2t$ in
Proposition~\ref{prop:AS-Poisson-gamma-morphism}.  Formula 26.4.21 identifies
the sum $u=0,\ldots,n-1$ with $Q(2t\mid2n)$, while the normal coordinate in
26.4.11 is
\[
 \frac{2t-2n}{\sqrt{4n}}=\frac{t-n}{\sqrt n}\longrightarrow0.
\]
Together with $P(0)=Q(0)=1/2$, this is the exact classical formula path to
Crandall's limit.  Formula 26.4.11 prints neither an error term nor a proof on
that page, however, so it is not used as an unspoken uniformity assertion.
The following argument controls all three real-$t$ endpoints directly.

Let $X_t$ have the Poisson distribution with mean $t$.  The characteristic
function of $Y_t=(X_t-t)/\sqrt t$ is
\[
 \mathbf E(e^{isY_t})
 =\exp\!\left(t\left(e^{is/\sqrt t}-1-is/\sqrt t\right)\right).
\]
The Taylor remainder is $O_s(t^{-1/2})$, so this converges to
$e^{-s^2/2}$.  The L\'evy continuity theorem gives
$Y_t\Rightarrow N(0,1)$.  Write $a_t=(k(t)-t)/\sqrt t$.  For every fixed
$\delta>0$, condition \eqref{eq:Crandall-Poisson-moving-cutoff} gives
\[
 \Pr(Y_t\leq-\delta)
 \leq \Pr(Y_t\leq a_t)
 \leq \Pr(Y_t\leq\delta)
\]
for all sufficiently large $t$.  Weak convergence at the two fixed continuity
points gives lower and upper limits between $\Phi(-\delta)$ and
$\Phi(\delta)$; letting $\delta\downarrow0$ proves
\eqref{eq:Crandall-Poisson-half-mass}.  Every cutoff in
\eqref{eq:Crandall-Poisson-three-cutoffs} differs from $t$ by at most a fixed
constant.  Removing the $u=0$ term changes the probability by exactly
$e^{-t}\to0$.

The high-resolution source page prints
$u<\lfloor t\rfloor$, not $u\leq\lfloor t\rfloor$.  Crandall points to
error-function asymptotics but prints no proof
\cite[p.~1287, Lem.~6.1]{Crandall1978}.  The argument above supplies the
all-real analytic completion while explicitly retaining its invocation of
L\'evy's continuity theorem.
\end{proof}

\begin{theorem}[Crandall's power-law theorem, with an explicit witness]
\label{thm:Crandall-power-law-count}
There are $c>0$ and $X_0>0$ such that
\begin{equation}
\label{eq:Crandall-power-law-count}
 \pi(x)>x^c\qquad(x>X_0).
\end{equation}
The source leaves $c$ unspecified.  Its proof, with the endpoint repaired,
permits the concrete edition choice $c=1/100$.
\end{theorem}

\begin{proof}
\phantomsection\label{sourceissue:Crandall-theorem61-endpoint}
For $x>x_0$, Theorem~\ref{thm:Crandall-fixed-height-count} may be summed only
over $h<r\log_2x$.  The printed proof sums through
$\lfloor r\log_2x\rfloor$, which includes one unauthorized equality term
when $r\log_2x$ is integral.  Summing instead through
$\lceil r\log_2x\rceil-1$ and applying
Lemma~\ref{lem:Crandall-Poisson-half-mass} gives, for every fixed
$d\in(0,1/2)$ and all sufficiently large $x$,
\begin{equation}
\label{eq:Crandall-power-intermediate}
 \pi(x)>(1/2-d)e^{r\log_2x}
       =(1/2-d)x^{r/\log2}.
\end{equation}
This proves the source theorem for every $c<r/\log2$ after increasing the
threshold.

For an exact explicit witness, take $r=1/100$.  Since
$3/64>2^{-5}$ and $\mathrm e<3$, the left side of
\eqref{eq:Crandall-r-choice} is greater than $45/100$, while its middle term
$3\mathrm e/100$ is less than $9/100$.  Thus this $r$ is admissible.  Since
$\log2<1$, one has $r/\log2>1/100$; the fixed positive factor in
\eqref{eq:Crandall-power-intermediate} is eventually absorbed by the excess
power.  Hence $c=1/100$ works.  This is an edition-supplied sharpening of the
source's existential conclusion, not a number printed by Crandall.
\end{proof}

\begin{nonclaim}
\label{nonclaim:Crandall-no-positive-density}
Theorem~\ref{thm:Crandall-power-law-count} is a lower bound of order $x^c$
for some $c<1$.  It is compatible with natural density zero and is not a
positive-density theorem.  It proves neither $\pi(x)\gg x$ nor convergence
for almost every odd starting value.  The frozen route text that calls this
positive density is routing metadata, not mathematical evidence.
\end{nonclaim}

\subsubsection{Membership coordinates and the return obstruction}

For a backward word $m=B_{b_k\cdots b_1}(n)$, Crandall uses the reverse
cumulative sums
\begin{equation}
\label{eq:Crandall-source-A-coordinates}
 A_0=0,
 \qquad
 A_i=\sum_{j=k-i+1}^{k}b_j\quad(1\leq i\leq k).
\end{equation}
Expanding the affine composition gives
\begin{equation}
\label{eq:Crandall-membership-identity}
 2^{A_k}n-3^km
 =\sum_{j=0}^{k-1}2^{A_j}3^{k-1-j}.
\end{equation}

\begin{theorem}[Crandall, exact trajectory-membership coordinates]
\label{thm:Crandall-membership}
For $m,n\in\Opos$, one has $n\in\mathcal T_m$ if and only if there are an
integer $k\geq1$ and integers
\[
 0=A_0<A_1<\cdots<A_k
\]
satisfying \eqref{eq:Crandall-membership-identity}.  If such data exist, then
$n=1$ or $n$ is the $k$-th displayed member of $\mathcal T_m$.
\end{theorem}

\begin{proof}
If $n=C_{\mathrm{Cr}}^k(m)$, reverse its $k$ branch equations and take
$b_i=e(C_{\mathrm{Cr}}^{k-i}(m))$; expansion gives
\eqref{eq:Crandall-membership-identity}.  Conversely, define
\begin{equation}
\label{eq:Crandall-recovered-differences}
 d_i=A_{k-i+1}-A_{k-i}\quad(1\leq i\leq k).
\end{equation}
The strict increase makes every $d_i$ positive, and solving
\eqref{eq:Crandall-membership-identity} gives
$m=B_{d_k\cdots d_1}(n)$.  Crandall's Lemma~4.1 then proves that every
formal inverse branch is integral, odd, and removed by the actual accelerated
map.  Hence $C_{\mathrm{Cr}}^k(m)=n$, with $n=1$ treated as the terminating
case \cite[pp.~1287--1288, Eq.~(7.2), Thm.~7.1]{Crandall1978}.
\end{proof}

\begin{sourceissue}[Theorem 7.1 difference index]
\label{sourceissue:Crandall-theorem71-index}
The printed converse assigns
$d_i=A_{k-i+1}-A_{k-1}$.  For $i=2$ this can give zero and it does not recover
successive coordinate gaps.  Equation~\eqref{eq:Crandall-recovered-differences}
is the forced correction.  It is exactly the difference needed to recover
$A_i$ from the $d$-word and makes the printed conclusion
$m=B_{d_k\cdots d_1}(n)$ valid.
\end{sourceissue}

Taking $n=m$ in \eqref{eq:Crandall-membership-identity} gives precisely
\eqref{eq:crandall-cycle-identity}.  The forward clause of Crandall's printed
Corollary~7.1 assumes that the trajectory has period $k$; its converse says
that the displayed equation produces a return at the $k$-th element.
Combining that converse with Theorem~\ref{thm:Crandall-membership} yields the
edition's slightly more general equivalence: the equation with strict
$A$-coordinates characterizes a return after $k$ accelerated steps.  It does
not assert that $k$ is the least period.  Positivity of the right side forces
\begin{equation}
\label{eq:Crandall-positive-cycle-denominator}
 2^{A_k}>3^k.
\end{equation}

\begin{proposition}[Crandall's explicit Diophantine obstruction]
\label{prop:Crandall-cycle-obstruction}
Suppose integers $a,b$ satisfy
\begin{equation}
\label{eq:Crandall-obstruction-hypothesis}
 b>1,
 \qquad
 0<2^a-3^b\mid 2^{a-b}-1.
\end{equation}
Put
\begin{equation}
\label{eq:Crandall-obstruction-construction}
 d=\frac{2^{a-b}-1}{2^a-3^b},
 \qquad
 m=2^bd-1.
\end{equation}
Then $m\in\Opos$, $m>1$, and $C_{\mathrm{Cr}}^b(m)=m$.  Therefore the
Collatz conjecture would be false.  The least period is allowed to divide
$b$.
\end{proposition}

\begin{proof}
The positive divisor relation forces $a>b$, so $d$ is a positive integer and
$m$ is positive odd.  If $m=1$, then $d=2^{1-b}$ is not an integer because
$b>1$; hence $m>1$.  Multiplying the definition of $d$ by $2^b$ gives
\[
 m(2^a-3^b)=3^b-2^b
 =\sum_{j=0}^{b-1}2^j3^{b-1-j}.
\]
This is the return equation with $k=b$, $A_j=j$ for $0\leq j<b$, and
$A_b=a$.  The strictness $A_{b-1}<A_b$ follows from $a>b$.  Corollary~
\ref{cor:crandall-cycle-identity} supplies the genuine return
\cite[p.~1288]{Crandall1978}.
\end{proof}

\subsubsection{Pillai's fixed-difference citation and the exact P\'olya route}

Crandall reports two external Diophantine results.  The first says that, for
each fixed integer $z$, the equation
\begin{equation}
\label{eq:Crandall-fixed-difference}
 2^x-3^y=z
\end{equation}
has only finitely many integer solutions $(x,y)$
\cite[ref.~9 and p.~1288]{Crandall1978}.  Pillai's cited primary paper has now
been read in full.  Its Corollary~1 states the corresponding positive-
exponent, positive-difference conclusion
\cite[p.~4]{Pillai1931}.  The paper also identifies the qualitative result as
a special case of P\'olya's earlier fixed-prime gap theorem
\cite[p.~2]{Pillai1931}.  Using P\'olya's sourced gap theorem yields exactly
the finiteness required by Crandall and bypasses Pillai's defective
quantitative derivation.  The proof-dependency audit of P\'olya's explicit
Thue citations is recorded separately below.

\begin{sourceissue}[Pillai's literal quantitative proof]
\label{sourceissue:Pillai-degree-gap}
Pillai quotes Siegel's approximation exponent correctly but calls its
constant a positive \emph{integer}; Siegel asserts only a positive real
constant \cite[p.~174]{Siegel1921}.  More seriously, Pillai puts
$\alpha=(A/B)^{1/r}$ and asserts that $\alpha$ has degree $r$ whenever $A$
and $B$ are not both perfect $r$th powers \cite[p.~3]{Pillai1931}.  This is
false: $(4/1)^{1/4}=\sqrt2$ has degree two, and
$(2/2)^{1/2}=1$ has degree one although neither displayed coefficient is a
square.  His next comparison bounds $Y/X$ only from above, although the step
which replaces $Y^{r-1}$ by a constant multiple of $X^{r-1}$ needs a lower
bound.  Finally, the proof of Theorem~I applies the lemma to every residue
pair without checking its ``not both powers'' hypothesis; the residue pair
$h=\ell=0$ with outer coefficients one violates it.  Thus the printed
derivation is not admitted literally.
\end{sourceissue}

The defect is repairable on the exact source theorem rather than by changing
Pillai's conclusion.  For $q\geq1$ write
\begin{equation}
\label{eq:Pillai-mq}
 \mu(q)=\min_{1\leq\lambda\leq q}
 \left(\frac{q}{\lambda+1}+\lambda\right).
\end{equation}

\begin{theorem}[Siegel's rational-approximation theorem]
\label{thm:Siegel-rational-approximation}
Let $\xi$ be algebraic of degree $q\geq2$ and let $\epsilon>0$.
There is a positive real constant $c=c(\xi,\epsilon)$ such that, for all
rational integers $u,v$ with $v>0$,
\[
 \left|\xi-\frac uv\right|
 >\frac{c}{v^{\mu(q)+\epsilon}}.
\]
The constant is not asserted to be an integer or numerically effective.
\end{theorem}

\begin{proof}
This is the denominator form printed on Siegel's p.~174
\cite[p.~174]{Siegel1921}.  His Hauptsatz is stated in height form on
printed p.~178 and proved through p.~191; Zus\"atze~1--3 give a uniform
positive constant and extend the conclusion to nonintegral algebraic
$\xi$ \cite[pp.~178, 191]{Siegel1921}.  Although the Hauptsatz minimizes over
$1\leq s<q$, this is the same minimum as in \eqref{eq:Pillai-mq}: the
excluded endpoint has value $q+q/(q+1)>q$, whereas the candidate $s=1$ has
value $q/2+1\leq q$.  If $u/v=p/s$ is reduced, then
$s\leq v$.  In the error-$<1$ case,
$H(p/s)=\max\{|p|,s\}\leq\max\{1,|\xi|+1\}s$; the complementary error
case is absorbed by decreasing the positive constant.  This is the exact
height-to-denominator map, with no assertion of an effective constant.
\end{proof}

\begin{lemma}[Monotonicity of Siegel's exponent]
\label{lem:Siegel-exponent-monotone}
The sequence $\mu(q)$ is nondecreasing on $\Npos$.
\end{lemma}

\begin{proof}
For an old candidate $1\leq\lambda\leq q$, replacing $q$ by $q+1$
increases $q/(\lambda+1)+\lambda$ by $1/(\lambda+1)$.  The only new
candidate is $\lambda=q+1$, whose value
$(q+1)/(q+2)+q+1$ is greater than $q+1$.  But
$\mu(q)\leq q/2+1\leq q+1$, using the old candidate $\lambda=1$.
Every candidate for $\mu(q+1)$ is therefore at least $\mu(q)$.
\end{proof}

\begin{proposition}[Actual-degree repair of Pillai's lemma]
\label{prop:Pillai-actual-degree-repair}
Fix $A,B,r\in\Npos$ with $r\geq2$ and fix $\eta>0$.  There is a constant
$C=C(A,B,r,\eta)>0$ such that, for all $X,Y\in\Npos$ with
$AX^r>BY^r$,
\begin{equation}
\label{eq:Pillai-repaired-power-separation}
 AX^r-BY^r>CX^{r-\mu(r)-\eta}.
\end{equation}
Consequently, for every $\delta>0$ the coefficient-free inequality
$AX^r-BY^r>X^{r-\mu(r)-\delta}$ holds for all sufficiently large $X$.
No hypothesis that $A$ and $B$ are not both $r$th powers is required.
\end{proposition}

\begin{proof}
Put $\alpha=(A/B)^{1/r}$ and
$d=[\Q(\alpha):\Q]\leq r$.  The exact factorization is
\begin{equation}
\label{eq:Pillai-exact-factorization}
 AX^r-BY^r
 =B(\alpha X-Y)
   \sum_{j=0}^{r-1}(\alpha X)^{r-1-j}Y^j.
\end{equation}
If $d\geq2$, Siegel's theorem at the \emph{actual} degree gives a constant
$c>0$ for which
\[
 \left|\alpha-\frac YX\right|>cX^{-\mu(d)-\eta}.
\]
The difference is positive, and Lemma~\ref{lem:Siegel-exponent-monotone}
gives $\mu(d)\leq\mu(r)$.  Hence
$\alpha X-Y>cX^{1-\mu(r)-\eta}$.  Retaining the positive $j=0$ term in
the sum in \eqref{eq:Pillai-exact-factorization} gives
\[
 AX^r-BY^r
 >Bc\alpha^{r-1}X^{r-\mu(r)-\eta}.
\]
If $d=1$, write $\alpha=P/Q$ in lowest positive terms.  Strict positivity
excludes $PX=QY$, so $PX-QY\geq1$ and $\alpha X-Y\geq1/Q$.  Since $r\geq2$
and $Y>0$, the factor sum in \eqref{eq:Pillai-exact-factorization} is strictly
larger than its $j=0$ term.  Hence
\[
 AX^r-BY^r>\frac{B\alpha^{r-1}}QX^{r-1}.
\]
Moreover $\mu(r)>1$, because every candidate in \eqref{eq:Pillai-mq} has
$\lambda\geq1$ and a strictly positive first summand.  The displayed bound
therefore implies \eqref{eq:Pillai-repaired-power-separation} for $X\geq1$.
Finally, with $\eta=\delta/2$, the ratio of the established lower bound to
$X^{r-\mu(r)-\delta}$ is $CX^{\delta/2}$, which exceeds one for
$X>C^{-2/\delta}$.
\end{proof}

\begin{theorem}[Repaired quantitative Pillai theorem]
\label{thm:Pillai-repaired-quantitative}
Fix $A,B,M,N\in\Npos$ with $M,N>1$ and $\delta>0$, and assume, as Pillai
does, that $\log M/\log N\notin\Q$.  Then there is an integer
$x_0=x_0(A,B,M,N,\delta)\geq0$ such that
\begin{equation}
\label{eq:Pillai-repaired-theorem-I}
 AM^x-BN^y>M^{x(1-\delta)}
\end{equation}
whenever $x,y\in\Npos$, $x>x_0$, and $AM^x>BN^y$.
\end{theorem}

\begin{proof}
Replace $\delta$ by $\delta_0=\min(\delta,1/2)$.  The required large $r$
exists constructively: for $k=\lceil\sqrt r\rceil$ and
$\lambda=k-1\geq1$,
\[
 \mu(r)\leq\frac r{k}+k-1<2\sqrt r,
\]
so $\mu(r)/r<2/\sqrt r$.  Choose, for example, an integer
$r>(4/\delta_0)^2$, increasing it to at least two if necessary.  Then
$\mu(r)/r<\delta_0/2$.  Choose $\eta>0$ with
$\eta/r<\delta_0/4$.  Write uniquely
\[
 x=rs+h,\qquad y=rt+\ell,
 \qquad 0\leq h,\ell<r.
\]
For the fixed residue pair $(h,\ell)$, put
$A_h=AM^h$, $B_\ell=BN^\ell$, $X=M^s$, and $Y=N^t$.
Proposition~\ref{prop:Pillai-actual-degree-repair} gives
\[
 AM^x-BN^y
 >C_{h,\ell}M^{(x-h)(1-(\mu(r)+\eta)/r)}.
\]
Put
\[
 C_*=\min_{0\leq h,\ell<r}C_{h,\ell}>0,
 \qquad
 \gamma=1-\frac{\mu(r)+\eta}{r},
 \qquad
 \varepsilon=\delta_0-\frac{\mu(r)+\eta}{r}>\frac{\delta_0}{4}.
\]
Here $\gamma>1-3\delta_0/4>0$.  For every residue pair,
\[
 \frac{C_{h,\ell}M^{(x-h)\gamma}}{M^{x(1-\delta_0)}}
 =C_{h,\ell}M^{x\varepsilon-h\gamma}
 \geq C_*M^{x\varepsilon-(r-1)\gamma}.
\]
The last quantity exceeds one whenever
\[
 x>\frac{(r-1)\gamma-\log_M C_*}{\varepsilon},
\]
a threshold independent of $y,h,\ell$.  This proves
$AM^x-BN^y>M^{x(1-\delta_0)}\geq M^{x(1-\delta)}$.
The preceding proof establishes the same quantified conclusion for all
$M,N>1$, without the irrational-log hypothesis.  This stronger statement is
an edition result; Pillai's sourced theorem remains recorded with his
hypothesis.
\end{proof}

We now type the independent qualitative route used for Crandall.  For a
finite ordered tuple $S=(p_1,\ldots,p_r)$ of pairwise distinct primes define
\[
 \mathcal A_S
 =\{p_1^{e_1}\cdots p_r^{e_r}:e_i\in\Nzero\}.
\]

\begin{proposition}[Prime-exponent coordinate morphism]
\label{prop:Polya-smooth-coordinate}
The map
\[
 E_S:(\Nzero^r,+)\longrightarrow(\mathcal A_S,\mathord\cdot),
 \qquad
 E_S(e_1,\ldots,e_r)=\prod_{i=1}^r p_i^{e_i},
\]
is a commutative-monoid isomorphism.  Its inverse is
$a\mapsto(\nu_{p_1}(a),\ldots,\nu_{p_r}(a))$; every fibre is a singleton.
Thus the coordinate loses no exponent information while the prime set $S$
is retained.
\end{proposition}

\begin{proof}
Multiplication in $\mathcal A_S$ adds every valuation, so $E_S$ preserves
the monoid law.  Existence and uniqueness of prime factorization prove both
surjectivity and the displayed two-sided inverse.
\end{proof}

\begin{theorem}[Thue's fixed-value theorem at P\'olya's exact scopes]
\label{thm:Thue-fixed-binary-form}
Let $c\in\Z\setminus\{0\}$.
\begin{enumerate}
\item If $U\in\Z[X,Y]$ is homogeneous and irreducible of degree greater
than two, then $U(x,y)=c$ has only finitely many solutions in $\Z^2$.
\item More generally, if a homogeneous integral binary form $F$ has
infinitely many integer solutions of $F(x,y)=c$, then, up to multiplication
by a nonzero scalar, $F$ is a power of a linear form or an indefinite
quadratic form.
\end{enumerate}
\end{theorem}

\begin{proof}
Thue's Theorem~IV is printed for positive $x,y$
\cite[pp.~303--304]{Thue1909}.  On the off-axis domain, the map
\[
 (x,y)\longmapsto
 \bigl(\operatorname{sgn}x,\operatorname{sgn}y,|x|,|y|\bigr)
\]
is a bijection
\[
 (\Z\setminus\{0\})^2
 \longrightarrow
 \{\pm1\}^2\times\Npos^2,
\]
with inverse $(\epsilon,\delta,X,Y)\mapsto(\epsilon X,\delta Y)$.
For each of the four sign pairs, the transformed form
$U_{\epsilon,\delta}(X,Y)=U(\epsilon X,\delta Y)$ remains integral,
homogeneous, irreducible, and of the same degree, so Thue's positive-quadrant
theorem gives a finite set.  If $y=0$, irreducibility in degree greater than
one forces the coefficient $u_d$ of $X^d$ to be nonzero, and the remaining
equation is $u_dx^d=c$; the other axis is identical.  This proves the exact
all-integer extension in part~1.

Part~2 is the exceptional-form formulation quoted by P\'olya
\cite[pp.~144--145]{Polya1918}.  His locator is Thue II, printed p.~3.  Both
available primary scans omit printed pp.~2--3, so that page is not represented
as visually verified.  Thue II printed p.~1 gives the equivalent
dehomogenized theorem and its degree-one/degree-two power exceptions
\cite[p.~1]{Thue1908II}.  The manifestation gap, P\'olya's exact quotation,
and the relation between the two formulations are recorded in
\path{qa/THUE-1908-1909-F017-binary-form-audit.md}.
\end{proof}

\begin{proposition}[P\'olya's residue-cell cubic morphism]
\label{prop:Polya-residue-cell-cubic}
Let
\[
 L_1(n)=an+b,\qquad L_2(n)=cn+d,
\]
where $a,b,c,d\in\Z$, $a,c\ne0$, and
$\Delta=ad-bc\ne0$.  Fix pairwise distinct primes
$p_1,\ldots,p_t$ and residue vectors
$\mathbf r,\mathbf s\in\{0,1,2\}^t$.  Let
$D_{\mathbf r,\mathbf s}$ be the set of integers $n$ for which
$L_1(n),L_2(n)>0$, both values have prime support in
$\{p_1,\ldots,p_t\}$, and
\[
 \nu_{p_i}(L_1(n))\equiv r_i\pmod3,\qquad
 \nu_{p_i}(L_2(n))\equiv s_i\pmod3
 \quad(1\leq i\leq t).
\]
Put
\[
 R=\prod_{i=1}^t p_i^{r_i},\qquad
 S_0=\prod_{i=1}^t p_i^{s_i},\qquad
 G_{\mathbf r,\mathbf s}(X,Y)=aS_0Y^3-cRX^3.
\]
Then
\[
 \Phi_{\mathbf r,\mathbf s}:D_{\mathbf r,\mathbf s}
 \longrightarrow
 \mathcal Z_{\mathbf r,\mathbf s}
 =\{(X,Y)\in\Npos^2:G_{\mathbf r,\mathbf s}(X,Y)=\Delta\},
\]
defined by
\[
 X=\prod_{i=1}^t
 p_i^{(\nu_{p_i}(L_1(n))-r_i)/3},
 \qquad
 Y=\prod_{i=1}^t
 p_i^{(\nu_{p_i}(L_2(n))-s_i)/3},
\]
is well defined and injective.  Its inverse on its image is
\[
 n=\frac{RX^3-b}{a}=\frac{S_0Y^3-d}{c}.
\]
The target $\mathcal Z_{\mathbf r,\mathbf s}$ is finite.
\end{proposition}

\begin{proof}
The residue hypotheses make the displayed exponents nonnegative integers and
give the exact identities
\[
 an+b=RX^3,\qquad cn+d=S_0Y^3.
\]
Multiplying the second identity by $a$, the first by $c$, and subtracting
gives
\[
 aS_0Y^3-cRX^3=ad-bc=\Delta.
\]
The two displayed recovery formulas agree by this equation and recover the
original $n$, proving injectivity and singleton fibres on the image.

The cubic has
\[
 \operatorname{Disc}(G_{\mathbf r,\mathbf s})
 =-27(aS_0)^2(cR)^2\ne0.
\]
It need not be irreducible over $\Q$, so part~1 of
Theorem~\ref{thm:Thue-fixed-binary-form} cannot be invoked alone.  Suppose
instead that, over a characteristic-zero field,
\[
 G_{\mathbf r,\mathbf s}(X,Y)=\lambda(uX+vY)^3.
\]
Its nonzero $X^3$ and $Y^3$ coefficients force
$\lambda,u,v\ne0$, whereas its zero $X^2Y$ coefficient would give
$3\lambda u^2v=0$, a contradiction.  It is therefore not a scalar linear
cube.  A quadratic power cannot have degree three.  Part~2 of
Theorem~\ref{thm:Thue-fixed-binary-form} now proves finiteness.  If one uses
P\'olya's positive-right-side phrasing and $\Delta<0$, the exact involution
\[
 (G_{\mathbf r,\mathbf s},\Delta,X,Y)
 \longmapsto
 (-G_{\mathbf r,\mathbf s},-\Delta,X,Y)
\]
preserves the solution set and exceptional-form status.
\end{proof}

\begin{theorem}[P\'olya's fixed-prime gap theorem]
\label{thm:Polya-smooth-gaps}
Let $r\geq2$ and enumerate $\mathcal A_S$ increasingly as
$a_0<a_1<a_2<\cdots$.  Then
\begin{equation}
\label{eq:Polya-smooth-gaps}
 \lim_{j\to\infty}(a_{j+1}-a_j)=\infty.
\end{equation}
\end{theorem}

\begin{proof}
This is the exact statement in P\'olya's Section~4, equations (12)--(14)
\cite[pp.~147--148]{Polya1918}.  P\'olya proves it equivalent to his Satz~I
for products of two essentially distinct rational linear factors.  His proof
of Satz~I reduces finitely supported prime-exponent vectors modulo three,
fixes a residue vector, and invokes Thue's finiteness theorem for the resulting
binary cubic \cite[pp.~144--145]{Polya1918}.  Proposition~
\ref{prop:Polya-residue-cell-cubic} verifies the exact dependency, including
the reducible-cubic cells not covered by Thue's irreducible-form theorem
alone.  Thue I, Satz~12, records the same modulo-three proof template
\cite[pp.~30--31]{Thue1908I}; P\'olya supplies the residue-cell reduction used
here.  No effective gap threshold is imported.
\end{proof}

\begin{corollary}[Finite fixed-difference fibres]
\label{cor:Polya-fixed-difference-fibres}
For every nonempty finite prime set $S$ and every fixed
$z\in\Z\setminus\{0\}$, the
fibre
\[
 \Delta_S^{-1}(z)
 =\{(u,v)\in\mathcal A_S^2:u-v=z\},
 \qquad \Delta_S(u,v)=u-v,
\]
is finite.
\end{corollary}

\begin{proof}
If $|S|\geq2$, choose an index after which every adjacent gap in
\eqref{eq:Polya-smooth-gaps} exceeds $|z|$.  The difference between any two
distinct later terms is a sum of one or more adjacent positive gaps, hence
also exceeds $|z|$.  Only pairs with at least one earlier coordinate remain.
If $a_0,\ldots,a_{J-1}$ are the earlier terms, then choosing one of these
values and which coordinate of $(u,v)$ it occupies leaves the other
coordinate uniquely determined by $u-v=z$; hence at most $2J$ pairs remain.
If
$S=\{p\}$, use directly $p^{j+1}-p^j=(p-1)p^j\to\infty$.
\end{proof}

\begin{proposition}[Exact fixed-$z$ consequence of P\'olya's sourced theorem]
\label{prop:Crandall-fixed-z-finiteness}
For every $z\in\Z$, equation \eqref{eq:Crandall-fixed-difference} has only
finitely many solutions $(x,y)\in\Z^2$.
\end{proposition}

\begin{proof}
First prove the exponent-domain reduction.  If $x=-a<0$ and $y\geq0$, then
\[
 2^x-3^y=\frac{1-2^a3^y}{2^a}
\]
has odd numerator and is not an integer.  If $x\geq0$ and $y=-b<0$, the
numerator of $(3^b2^x-1)/3^b$ is $-1$ modulo $3$.  If both exponents are
negative, then
\[
 2^{-a}-3^{-b}=\frac{3^b-2^a}{2^a3^b}
\]
has absolute value strictly between zero and one: its numerator is nonzero
by unique factorization and has absolute value smaller than the denominator.
Thus an integer-valued difference forces $x,y\in\Nzero$.

Define
\[
 I:\Nzero^2\longrightarrow\mathcal A_{\{2,3\}}^2,
 \qquad I(x,y)=(2^x,3^y).
\]
The valuation map
\[
 V(u,v)=(\nu_2(u),\nu_3(v))
\]
satisfies $V\circ I=\operatorname{id}_{\Nzero^2}$, so $I$ is injective.  For
$z\neq0$,
the solution set is exactly
\[
 I^{-1}\!\left(\Delta_{\{2,3\}}^{-1}(z)\right).
\]
Corollary~\ref{cor:Polya-fixed-difference-fibres} makes the inner fibre
finite, and injectivity makes its preimage finite.  For $z=0$, unique
factorization gives only
$2^x=3^y=1$, hence $(x,y)=(0,0)$.  This proves Crandall's literal integer-
exponent statement; it is not merely the positive-$z$ statement printed in
Pillai's Corollary~1.
\end{proof}

\subsubsection{Herschfeld's eventual-uniqueness theorem}

Crandall's reference~10 spells the author's surname ``Herschefeld.''  The
primary paper's title and byline give \emph{Aaron Herschfeld}
\cite[pp.~231--234]{Herschfeld1936}.  Crandall reports at most one integer
solution for sufficiently large $z$ \cite[p.~1288]{Crandall1978}.  The
primary theorem is stronger in the sign parameter but narrower in the
exponent domain: it treats sufficiently large $|d|$ and positive exponents.
The exact domain bridge is proved below rather than hidden in the citation.

\begin{sourceissue}[Herschfeld source boundaries]
\label{sourceissue:Herschfeld-boundaries}
On printed p.~232, a congruence argument for the specialized equation (2)
ends by referring to equation (1); the objects and displayed congruence make
the intended label unambiguous, but the defect is retained.  Printed p.~233
reports a further $|d|\leq100$ calculation from an unpublished paper; it is
not used here.  The final paragraph's general at-most-nine result for
$a^x-b^y=d$ has a separate Siegel-1929 dependency and is not a premise of the
special $(2,3)$ theorem.  That general result remains outside the admitted
scope until its exact cited theorem and coefficient hypotheses are read.
\end{sourceissue}

\begin{lemma}[The two order coordinates]
\label{lem:Herschfeld-orders}
For $r\geq1$ and $s\geq3$,
\begin{equation}
\label{eq:Herschfeld-orders}
 \operatorname{ord}_{3^r}(2)=2\cdot3^{r-1},
 \qquad
 \operatorname{ord}_{2^s}(3)=2^{s-2}.
\end{equation}
\end{lemma}

\begin{proof}
If $2^n\equiv1\pmod{3^r}$, reduction modulo three shows that $n=2m$ is
even.  The lifting-the-exponent identity gives
\[
 v_3(2^{2m}-1)=v_3(4^m-1)=1+v_3(m).
\]
Thus $3^r\mid2^n-1$ exactly when $2\cdot3^{r-1}\mid n$.  If
$3^n\equiv1\pmod{2^s}$ with $s\geq3$, reduction modulo eight shows that
$n$ is even, and
\[
 v_2(3^n-1)=v_2(3-1)+v_2(3+1)+v_2(n)-1=2+v_2(n).
\]
Therefore $2^s\mid3^n-1$ exactly when $2^{s-2}\mid n$.  This also makes
precise Herschfeld's older phrase that $3$ ``belongs to'' $2^{s-2}$ modulo
$2^s$ \cite[pp.~232--234]{Herschfeld1936}.
\end{proof}

\begin{theorem}[Herschfeld's positive-exponent theorem, with dependency repaired]
\label{thm:Herschfeld-positive-eventual-uniqueness}
There is a real number $D>0$ such that, for every $d\in\Z$ with $|d|>D$,
\begin{equation}
\label{eq:Herschfeld-positive-fibre}
 \#\{(x,y)\in\Npos^2:2^x-3^y=d\}\leq1.
\end{equation}
The proof supplies no effective numerical value of $D$.
\end{theorem}

\begin{proof}
Fix $0<\delta<1/2$.  Apply Theorem~
\ref{thm:Pillai-repaired-quantitative} first to bases $(2,3)$ and then to
$(3,2)$.  Let $x_1,y_1\in\Nzero$ be thresholds such that
\begin{align}
 0<2^u-3^v,\quad u>x_1
 &\Longrightarrow 2^u-3^v>2^{u(1-\delta)},
 \label{eq:Herschfeld-Pillai-positive}\\
 0<3^v-2^u,\quad v>y_1
 &\Longrightarrow 3^v-2^u>3^{v(1-\delta)}.
 \label{eq:Herschfeld-Pillai-negative}
\end{align}
Set
\begin{equation}
\label{eq:Herschfeld-threshold}
 D=\max\{2^{x_1+5},3^{y_1+2}\}.
\end{equation}

Suppose two distinct solutions have the same $d$.  They cannot agree in one
coordinate without agreeing in both.  Relabel them so that $X>x$; then
\[
 2^X-2^x=3^Y-3^y>0,
\]
so $Y>y$.  Factoring and using coprimality gives
\begin{equation}
\label{eq:Herschfeld-two-solution-factorization}
 2^x(2^{X-x}-1)=3^y(3^{Y-y}-1),
\end{equation}
and hence
\begin{equation}
\label{eq:Herschfeld-two-orders}
 2^{X-x}\equiv1\pmod{3^y},
 \qquad
 3^{Y-y}\equiv1\pmod{2^x}.
\end{equation}
Lemma~\ref{lem:Herschfeld-orders} yields
\begin{equation}
\label{eq:Herschfeld-exponent-separation}
 X-x\geq2\cdot3^{y-1},
 \qquad
 Y-y\geq2^{x-2}\quad(x>2).
\end{equation}

If $d>0$, then $d>2^{x_1+5}$ and
$d=2^x-3^y<2^x$, so $x>x_1+5>5$.  Therefore
\[
 2^X=3^Y+d>3^Y
 \geq3^{2^{x-2}}>2^{2^{x-2}},
 \qquad X>2^{x-2}>2x.
\]
Applying \eqref{eq:Herschfeld-Pillai-positive} at the larger solution gives
\[
 d=2^X-3^Y>2^{X(1-\delta)}>2^{X/2}>2^x>d,
\]
a contradiction.

If $d<0$, then $|d|>3^{y_1+2}$ and
$|d|=3^y-2^x<3^y$, so $y>y_1+2>2$.  Now
\[
 3^Y=2^X-d>2^X
 \geq2^{2\cdot3^{y-1}}>3^{3^{y-1}},
 \qquad Y>3^{y-1}>2y.
\]
Applying \eqref{eq:Herschfeld-Pillai-negative} at the larger solution gives
\[
 |d|=3^Y-2^X>3^{Y(1-\delta)}>3^{Y/2}>3^y>|d|,
\]
again a contradiction.  The elementary inequalities used above are strict
for exactly the displayed ranges: $2^{x-2}>2x$ for $x>5$ and
$3^{y-1}>2y$ for $y>2$.  This reconstructs Herschfeld's printed proof while
replacing only its reliance on Pillai's defective literal derivation by the
separately attributed repair already proved in this edition
\cite[pp.~233--234]{Herschfeld1936}.
\end{proof}

The next statement is an edition consequence, not part of Herschfeld's
printed exponent domain.

\begin{proposition}[Exact all-integer extension]
\label{prop:Herschfeld-integer-eventual-uniqueness}
For the same $D$ as in \eqref{eq:Herschfeld-threshold} and every
$d\in\Z$ with $|d|>D$,
\begin{equation}
\label{eq:Herschfeld-integer-fibre}
 \#\{(x,y)\in\Z^2:2^x-3^y=d\}\leq1.
\end{equation}
\end{proposition}

\begin{proof}
The exponent-domain argument in Proposition~
\ref{prop:Crandall-fixed-z-finiteness} proves that any integer-valued
difference has $x,y\in\Nzero$.  It remains to include a possible zero
coordinate without importing it into Herschfeld's statement.

Suppose two nonnegative solutions are ordered as in the preceding proof.  If
$d>0$, then $d<2^x$ still forces $x>x_1+5>5$, even if $y=0$.  The second
congruence in \eqref{eq:Herschfeld-two-orders} and the second bound in
\eqref{eq:Herschfeld-exponent-separation} remain valid, and the larger pair
$(X,Y)$ has positive coordinates.  The positive-$d$ contradiction is
unchanged.  If $d<0$, then $|d|<3^y$ still forces $y>y_1+2>2$, even if
$x=0$.  The first congruence and first order bound remain valid, the larger
pair is positive, and the negative-$d$ contradiction is unchanged.  Thus
the nonnegative fibre has size at most one.  Composing this conclusion with
the proved integer-to-nonnegative exponent reduction gives
\eqref{eq:Herschfeld-integer-fibre}.  This proves the exact domain used by
Crandall and strengthens his one-sided large-$z$ wording to Herschfeld's
two-sided large-$|d|$ parameter without changing either source statement.
\end{proof}

\begin{nonclaim}
Herschfeld's eventual uniqueness is not inferred from P\'olya's fixed-$d$
finiteness, and fixed-$d$ finiteness is not inferred from eventual
uniqueness.  The threshold in \eqref{eq:Herschfeld-threshold} is ineffective;
the theorem supplies neither an effective search bound nor a Collatz
endpoint or cycle exclusion.
\end{nonclaim}

\begin{proposition}[Exact primitive-root solution-to-return construction]
\label{prop:Crandall-primitive-root-return}
Put $t=\log_2 3$.  Let $p$ be an odd prime such that
\(\operatorname{ord}_p(2)=p-1\), and suppose
\begin{equation}
\label{eq:Crandall-primitive-root-input}
 (x,y)\in\Z^2,
 \qquad 2^x-3^y=p,
 \qquad y\geq\frac{p}{t-1}.
\end{equation}
Then one can construct integers
\[
 0=A_0<A_1<\cdots<A_y=x
\]
and a positive odd integer $m>1$ such that
\begin{equation}
\label{eq:Crandall-primitive-root-return}
 m(2^{A_y}-3^y)
 =\sum_{j=0}^{y-1}2^{A_j}3^{y-1-j}.
\end{equation}
Consequently $m$ occurs in its own Crandall trajectory, at accelerated
return time $y$; its least period is only asserted to divide $y$.
\end{proposition}

\begin{proof}
The integer-exponent denominator argument in the proof of Proposition~
\ref{prop:Crandall-fixed-z-finiteness} gives $x,y\in\Nzero$.  Under the last
hypothesis in \eqref{eq:Crandall-primitive-root-input}, $y>0$.  The equality
then excludes $p=3$ after reduction modulo $3$, so $p\geq5$.  Since
$0<t-1<1$, the last hypothesis also gives $y>p\geq5$; in particular all
index ranges below are nonempty.

Because $2^x>3^y$, one has $x>ty$.  Hence the integer
\begin{equation}
\label{eq:Crandall-primitive-root-slack}
 L=x-y>(t-1)y\geq p
 \qquad\Longrightarrow\qquad L\geq p+1.
\end{equation}
Equivalently, the assumed lower bound on $y$ is the exact integer inequality
$3^y\geq2^{p+y}$; no decimal approximation to $t$ is being used.

First put $A_j=j$ for $0\leq j\leq y-2$, and set
\[
 R_0=\sum_{j=0}^{y-2}2^j3^{y-1-j}.
\]
If $R_0\not\equiv0\pmod p$, choose $A_{y-1}$ to be the least integer in
$[y-1,x-1]$ satisfying
\begin{equation}
\label{eq:Crandall-primitive-root-selector-first}
 2^{A_{y-1}}\equiv-R_0\pmod p.
\end{equation}
Such an exponent exists: the interval contains at least $p-1$ consecutive
integers by \eqref{eq:Crandall-primitive-root-slack}, and
\(\operatorname{ord}_p(2)=p-1\) says exactly that the powers of $2$ over any
complete exponent-residue system modulo $p-1$ biject onto
\((\Z/p\Z)^\times\).

If instead $R_0\equiv0\pmod p$, keep $A_j=j$ for $0\leq j\leq y-3$, replace
$A_{y-2}$ by $y-1$, and write
\[
 R_1=\sum_{j=0}^{y-3}2^j3^{y-1-j}+2^{y-1}3.
\]
The replacement changes the partial sum by
\[
 R_1-R_0=3\cdot2^{y-2}\not\equiv0\pmod p,
\]
so $R_1\not\equiv0\pmod p$.  Choose $A_{y-1}$ to be the least integer in
$[y,x-1]$ satisfying
\begin{equation}
\label{eq:Crandall-primitive-root-selector-second}
 2^{A_{y-1}}\equiv-R_1\pmod p.
\end{equation}
This interval again contains at least $p-1$ consecutive integers, since its
cardinality is $L\geq p+1$.  Thus the selector is defined in both cases, and
the resulting sequence is strictly increasing with $A_0=0$ and $A_y=x$.

Let
\[
 S=\sum_{j=0}^{y-1}2^{A_j}3^{y-1-j}.
\]
The selected congruence gives $p\mid S$.  The $j=0$ summand is odd and every
other summand is even, so $S$ and $m=S/p$ are odd.  Moreover $0<t-1<1$
gives $y>p$, whence
\[
 S\geq3^{y-1}\geq3^p>p;
\]
The last strict inequality follows by induction from $3^n>n$ for $n\geq1$;
therefore $m$ is a positive odd integer greater than $1$.  Finally,
$2^{A_y}-3^y=2^x-3^y=p$, so $mp=S$ is exactly
\eqref{eq:Crandall-primitive-root-return}.  Corollary~7.1, in its exact form
proved above, now gives the asserted self-return.
\end{proof}

\begin{corollary}[Crandall's conditional primitive-root bound]
\label{cor:Crandall-primitive-root-bound}
Assume Conjecture~2.1.  If $p$ is a prime for which $2$ is a primitive root,
then every integer solution of $2^x-3^y=p$ satisfies
\[
 y<\frac{p}{\log_2 3-1}.
\]
\end{corollary}

\begin{proof}
A solution violating the strict bound satisfies the hypotheses of
Proposition~\ref{prop:Crandall-primitive-root-return} and therefore produces
a positive odd $m>1$ of infinite height, contrary to Conjecture~2.1.
\end{proof}

\begin{nonclaim}
\label{nonclaim:Crandall-primitive-root}
The source prints the conditional bound but omits this construction.  The
proposition does not assert that any solution violating the bound exists, and
it is not the adjacent sufficient obstruction with $A_j=j$.  The selected
return time $y$ need not be the least period, and the corollary remains
conditional on the Collatz conjecture.
\end{nonclaim}

\subsubsection{Continued fractions and the repaired cycle bound}

Put $t=\log_2 3$, and let $p_n/q_n$ be its continued-fraction convergents,
with
\begin{equation}
\label{eq:Crandall-convergent-recurrence}
 p_{-1}=1,\ q_{-1}=0,\ p_0=a_0,\ q_0=1,
 \qquad
 p_n=a_np_{n-1}+p_{n-2},\quad
 q_n=a_nq_{n-1}+q_{n-2}.
\end{equation}

\begin{lemma}[Exact continued-fraction transfer for the Crandall variables]
\label{lem:Crandall-CF-dependency}
The number $t=\log_2 3$ is irrational.  Let $p_n/q_n$ be its simple
continued-fraction convergents, with the recurrence
\eqref{eq:Crandall-convergent-recurrence}.  Then, for $n\geq2$,
\begin{equation}
\label{eq:Crandall-CF-error-bounds}
 \frac{1}{q_n+q_{n+1}}
 < |p_n-q_nt|
 < \frac{1}{q_{n+1}}.
\end{equation}
If $x\in\Z$ and $0<y<q_n$, then
\begin{equation}
\label{eq:Crandall-CF-best-approximation}
 |p_n-q_nt|<|x-yt|.
\end{equation}
For $n\geq4$, if $y=q_n$ and $x\in\Z$ with $x\ne p_n$, then the same
strict inequality holds with $y=q_n$.
\end{lemma}

\begin{proof}
Since $t>1$, a rational representation $t=a/b$ with $b>0$ would have
$a>0$ and $2^a=3^b$, contradicting unique factorization.  The continued
fraction is therefore infinite.  Steuding's Theorems~3.1--3.2
\cite[printed p.~38]{Steuding2005} give the
displayed recurrence and
\[
 p_nq_{n-1}-p_{n-1}q_n=(-1)^{n-1},
\]
so each $p_n/q_n$ is reduced.  Steuding's Theorem~3.6
\cite[printed p.~42]{Steuding2005}, applied to the
complete tail $\alpha_{n+1}$ of $t$, gives
\[
 t-\frac{p_n}{q_n}
 =\frac{(-1)^n}{q_n(\alpha_{n+1}q_n+q_{n-1})}.
\]
Because $a_{n+1}<\alpha_{n+1}<a_{n+1}+1$ and
$q_{n+1}=a_{n+1}q_n+q_{n-1}$, multiplying by $q_n$ yields
\eqref{eq:Crandall-CF-error-bounds}.  For $x>0$, Steuding's
Theorem~3.8 \cite[printed pp.~44--45]{Steuding2005}, with $(p,q)=(x,y)$,
gives \eqref{eq:Crandall-CF-best-approximation};
the equality case $x/y=p_n/q_n$ cannot occur because $p_n/q_n$ is reduced
and $0<y<q_n$.  If $x\leq0$, then $|x-yt|\geq yt>1$, since
$t=\log_2 3>1$, while the upper bound
in \eqref{eq:Crandall-CF-error-bounds} is $<1$.

For the final assertion, \eqref{eq:Crandall-CF-error-bounds} gives
$|p_n-q_nt|<1/q_{n+1}<1/2$ for $n\geq4$.  If $x\ne p_n$, integrality gives
$|x-p_n|\geq1$, and the triangle inequality gives
$|x-q_nt|\geq1-|p_n-q_nt|>|p_n-q_nt|$.
\end{proof}

\begin{sourceissue}[four repairs in the continued-fraction chain]
\label{sourceissue:Crandall-continued-fraction-chain}
The proof chain on printed pp.~1289--1290 needs four distinct repairs
\cite[pp.~1289--1290]{Crandall1978}.
\begin{enumerate}[label=(\roman*)]
\item Lemma~7.1 omits $0<y$.  As printed, $(x,y)=(0,0)$ contradicts its
strict best-approximation inequality.  The used form has $0<y<q_n$ and is
proved from Lemma~\ref{lem:Crandall-CF-dependency}.
\item Lemma~7.3 introduces $p_n/q_n$ but then writes $p_m,q_m$, leaves $x$
unquantified, and uses $|e^u-1|>|u|$.  This analytic inequality holds for
$u>0$, not for $u<0$.  The exact form needed for cycles is
\begin{equation}
\label{eq:Crandall-repaired-exponential-bound}
 \left.
 \begin{gathered}
  x,y\in\Z,\quad 0<y<q_n,\\
  x-y\log_2 3>0
 \end{gathered}
 \right\}
 \Longrightarrow
 |2^x-3^y|>3^y\log2\,|p_n-q_nt|.
\end{equation}
Indeed, for $u=(x-yt)\log2>0$ the positive-term power series gives
$e^u-1-u=\sum_{j\geq2}u^j/j!>0$, and best approximation supplies
$|x-yt|>|p_n-q_nt|$.  In the cycle application, the missing sign is exactly
\eqref{eq:Crandall-positive-cycle-denominator}.
\item Theorem~7.2 explicitly reduces to $k\leq q_n$, but then invokes
Lemma~7.3, whose repaired denominator hypothesis is strict: $y<q_n$.
The final assertion of Lemma~\ref{lem:Crandall-CF-dependency} handles
$k=q_n$ when $A_k\ne p_n$.  If $A_k=p_n$, the positive exponential
inequality is itself strict.  Thus the same denominator estimate holds at
the equality boundary that the printed proof does not justify.
\item Corollary~7.2 claims that its minimum becomes unbounded as $n$
increases.  For a fixed cycle its second entry instead tends to zero.  The
correct proof first fixes $n$ and only then chooses a sufficiently large
cycle minimum; this quantifier order is proved below.
\end{enumerate}
The later references to Eq.~(7.5) when computing $q_n$ are correct only as
references to the printed continued-fraction list; the recurrence itself is
Eq.~(7.6).
\end{sourceissue}

\begin{theorem}[Crandall's cycle bound, with proof gaps repaired]
\label{thm:Crandall-cycle-bound}
Let $m>1$ be the least member of a $C_{\mathrm{Cr}}$-cycle of least
accelerated period $k$.  For every integer $n\geq4$,
\begin{equation}
\label{eq:Crandall-cycle-bound}
 k>\min\left\{q_n,\frac{2m}{q_n+q_{n+1}}\right\}.
\end{equation}
Consequently, only finitely many cycles have any prescribed accelerated
period \cite[p.~1290]{Crandall1978}.
\end{theorem}

\begin{proof}
Minimality of $m$ and the branch equation give, for the coordinates
\eqref{eq:Crandall-cumulative-time},
\begin{equation}
\label{eq:Crandall-cycle-growth-bound}
 2^{A_j}\leq(3+1/m)^j.
\end{equation}
The return equation and this bound imply
\begin{equation}
\label{eq:Crandall-cycle-minimum-bound}
 m<\frac{k(3+1/m)^{k-1}}{2^{A_k}-3^k}.
\end{equation}
If $k>q_n$, \eqref{eq:Crandall-cycle-bound} is immediate.  Otherwise the
repaired best-approximation argument in
Source issue~\ref{sourceissue:Crandall-continued-fraction-chain}, including
its equality case, gives
\[
 2^{A_k}-3^k>3^k\log2\,|p_n-q_nt|.
\]
Combining this with \eqref{eq:Crandall-cycle-minimum-bound}, and writing
$v=1/(3m)>0$, first yields
\[
 k>m(3+1/m)(\log2)|p_n-q_nt|(1+v)^{-k}.
\]
The function $f(v)=(1+v)^{-k}$ has
$f''(v)=k(k+1)(1+v)^{-k-2}>0$; strict convexity above its tangent at zero
therefore gives $f(v)>f(0)+f'(0)v=1-kv$.  Hence
\begin{equation}
\label{eq:Crandall-cycle-bound-intermediate}
 k>m(\log2)(3+1/m)|p_n-q_nt|\left(1-\frac{k}{3m}\right).
\end{equation}
For $n\geq4$, the recurrence gives
$q_n+q_{n+1}\geq q_4+q_5=12+41=53>20$.  If $k\geq m/10$, the second member
of the minimum in \eqref{eq:Crandall-cycle-bound} is strictly below
$m/10\leq k$.  Otherwise $k<m/10$, so
$1-k/(3m)>29/30$, and the continued-fraction lower bound
$|p_n-q_nt|>(q_n+q_{n+1})^{-1}$ turns
\eqref{eq:Crandall-cycle-bound-intermediate} into
\[
 k>\frac{(3m+1)\log2}{q_n+q_{n+1}}\frac{29}{30}
   >\frac{2m}{q_n+q_{n+1}}.
\]
The final strict inequality is exact: the positive-term expansion
$\log2=2(1/3+1/(3^3\cdot3)+\cdots)$ gives
$\log2>56/81>20/29$, while $(3m+1)/m>3$.
This proves \eqref{eq:Crandall-cycle-bound}.

For finiteness, fix a proposed period $L$ and choose an integer $n\geq4$ with
$q_n>L$; the convergent denominators are unbounded.  An infinite family of
distinct period-$L$ cycles would have unbounded distinct
least members, so choose one with
$m>L(q_n+q_{n+1})/2$.  Both entries of the minimum would then exceed $L$,
contradicting \eqref{eq:Crandall-cycle-bound} with $k=L$.
\end{proof}

\phantomsection\label{historical-10^9-premise}
Printed p.~1282 says that Conjecture~(2.1) had been verified for
$m<10^9$, citing references~[1] and~[5] \cite[p.~1282]{Crandall1978}.
Those references are not two independent witnesses.  Reference~[1] is
Conway's \emph{Unpredictable iterations}; its correct proceedings pages are
49--52, not Crandall's 45--52.  Reference~[5], printed as
``\emph{Zentralblatt f\"ur Mathematik}, Band~233, p.~10041,'' is a defective
rendering of \emph{Zbl}~0337.10041, the Zentralblatt record for that same
Conway paper: $10041$ is part of the accession identifier, not a printed
page.  Thus the two routes collapse to one work, rather than corroborating
one another \cite[p.~1292]{Crandall1978}.

Conway defines the ordinary one-step map $g$ by $g(n)=n/2$ for even $n$ and
$g(n)=3n+1$ for odd $n$, literally the map $U$ in
\eqref{eq:unshortened-U}.  On the first page of his article he reports the
inclusive range $n\leq10^9$ and attributes its verification to D.~H. Lehmer,
Emma Lehmer, and J.~L. Selfridge \cite[p.~219]{Conway1972}.  He supplies no
citation for that sentence and no program, machine, range partition,
stopping criterion, independent check, checksum, log, or preserved output.
His proved undecidability theorem is expressly said to say nothing about the
particular Collatz game; the verification sentence is historical background,
not a theorem of the article.  The published contemporaneous report and its
attribution are therefore content-verified, while the underlying computation
remains unauthenticated and nonreproducible at this checkpoint.

\begin{proposition}[One-step verification to accelerated exclusion]
\label{prop:Conway-Crandall-historical-transfer}
Let $N\in\Npos$.  Suppose that for every $n\in\Npos$ with $n\leq N$ there
exists $r\in\Nzero$ such that $U^r(n)=1$.  Then for every
$m\in\Opos$ with $m\leq N$ there exists $j\in\Nzero$ such that
$C_{\mathrm{Cr}}^j(m)=1$.  Consequently no $m\in\Opos$ with
$1<m\leq N$ can satisfy $C_{\mathrm{Cr}}^k(m)=m$ for any $k\in\Npos$.
\end{proposition}

\begin{proof}
Starting from an odd $m$, the successive odd states of the $U$-orbit are
exactly the $C_{\mathrm{Cr}}$-orbit, with the $j$th odd return occurring at
the unshortened time $A_j+j$ by \eqref{eq:Crandall-three-clocks}.  Since $1$
is odd, any occurrence $U^r(m)=1$ is one of those successive odd returns;
hence $C_{\mathrm{Cr}}^j(m)=1$ for a unique first such index $j$.  Finally
$C_{\mathrm{Cr}}(1)=1$.  If $m>1$ also satisfied
$C_{\mathrm{Cr}}^k(m)=m$, its accelerated orbit would be periodic from $m$
and could not enter the fixed point $1$, a contradiction.
\end{proof}

Printed p.~1290 uses the ``known result'' that no $1<m<10^9$ appears in
its own trajectory, without a new inline citation
\cite[p.~1290]{Crandall1978}.  Proposition
\ref{prop:Conway-Crandall-historical-transfer} proves the exact clock and
domain transfer from Conway's stronger inclusive ordinary-map report to that
accelerated statement; it does not authenticate the computation.

For a hypothetical nontrivial cycle, its least member $m_0$ lies in
$\Opos$.  Conway's inclusive report, after the proved transfer, directly
excludes $m_0\leq10^9$.  Independently, Crandall's weaker strict-range form
first gives $m_0\geq10^9$, and oddness excludes equality with the even
endpoint.  Both routes therefore give, exactly as Crandall writes,
\[
 m_0>10^9,
 \qquad\text{equivalently}\qquad m_0\geq1{,}000{,}000{,}001.
\]
With
\begin{equation}
\label{eq:Crandall-historical-convergents}
 q_{10}=31867,
 \qquad q_{11}=79335,
\end{equation}
Theorem~\ref{thm:Crandall-cycle-bound} gives the least accelerated period
$\ell$ the bound
\[
 \ell>\min\left\{31867,\frac{2\cdot10^9}{111202}\right\}>17985.
\]
Every positive return exponent $k$ is a multiple of $\ell$, so this recovers
Crandall's Theorem~7.3: $C_{\mathrm{Cr}}^k(m)=m$ with $m>1$ implies
$k>17985$, conditional on the published 1972 report reused by Crandall in
1978.  The source variable is an accelerated return time, not a current or
reproducible verification record.
For the same
return, \eqref{eq:Crandall-three-clocks} gives
$A_k\geq k>17985$ shortened steps and $A_k+k>k>17985$ unshortened steps;
these are transferred inequalities, not an identification of the three
clocks.

\phantomsection\label{sourceissue:Crandall-bounded-supremum}
The opening of Section~7 also says every member of a bounded trajectory is
strictly below its supremum.  The correct relation is $\leq$.  The intended
finite pigeonhole argument survives: boundedness plus absence of nontrivial
self-returns would force the trajectory to contain $1$.

\subsubsection{The generalized \texorpdfstring{$qx+r$}{qx+r} programme}

For positive odd $q,r$, $q>1$, Crandall defines
\begin{equation}
\label{eq:Crandall-qxr}
 C_{q,r}(m)=\frac{qm+r}{2^{e_{q,r}(m)}},
 \qquad
 e_{q,r}(m)=\vtwo(qm+r),
 \qquad m\in\Opos.
\end{equation}
Conjecture~8.1 predicts that, except for $(q,r)=(3,1)$, some positive odd
starting value fails to reach $1$.  The exceptional pair is not an assertion
that the ordinary conjecture is true.

\begin{proposition}[Crandall's proved generalized failures]
\label{prop:Crandall-generalized-failures}
The following statements hold.
\begin{enumerate}[label=(\alph*)]
\item If $r>1$, every $m\in\Opos$ divisible by $r$ fails to reach $1$.
\item For $(q,r)=(5,1)$,
\[
 13\longmapsto33\longmapsto83\longmapsto13,
 \qquad(e_0,e_1,e_2)=(1,1,5).
\]
\item For $(q,r)=(181,1)$,
\[
 27\longmapsto611\longmapsto27,
 \qquad(e_0,e_1)=(3,12).
\]
\end{enumerate}
\end{proposition}

\begin{proof}
If $r\mid m$, then $qm+r\equiv0\pmod r$; division by a power of $2$, a unit
modulo odd $r$, preserves divisibility by $r$.  Induction excludes the value
$1$.  The two cycles follow by exact substitution:
\[
 5(13)+1=2(33),\quad5(33)+1=2(83),\quad5(83)+1=2^5(13),
\]
and
\[
 181(27)+1=2^3(611),\qquad181(611)+1=2^{12}(27).
\]
The associated differences are $2^7-5^3=3$ and
$2^{15}-181^2=7$.  The printed starting value in the second cycle is $27$;
$21$ is an OCR error \cite[p.~1291]{Crandall1978}.
\end{proof}

\begin{proposition}[the omitted finite certificate for $q=1093$]
\label{prop:Crandall-q1093}
For $C_{1093,1}$, a positive odd integer has height one exactly when
\begin{equation}
\label{eq:Crandall-q1093-height-one}
 m=\frac{2^{364p}-1}{1093}
 \qquad(p\in\Z_{>0}).
\end{equation}
No positive odd integer has height two.  Every positive odd integer not in
\eqref{eq:Crandall-q1093-height-one} therefore has infinite height, and this
infinite-height set has natural density one in $\Opos$.
\end{proposition}

\begin{proof}
Exact modular exponentiation gives
\begin{equation}
\label{eq:Crandall-q1093-certificate}
 \operatorname{ord}_{1093}(2)=364,
 \qquad
 2^{364}\equiv1\pmod{1093^2}.
\end{equation}
For the first equality it suffices to check the proper divisors of
$364=2^2\cdot7\cdot13$; convenient witnesses include
\[
 2^{182}\equiv1092,
 \qquad2^{52}\equiv27,
 \qquad2^{28}\equiv121\pmod{1093}.
\]
A height-one value satisfies $1093m+1=2^a$, so the order calculation is
equivalent to $364\mid a$ and proves
\eqref{eq:Crandall-q1093-height-one}.  The congruence modulo $1093^2$ shows
that every such $m$ is itself divisible by $1093$.  If it had a predecessor
$u$, then $1093u+1=2^b m$ would be simultaneously congruent to $1$ and $0$
modulo $1093$, impossible.  Thus height two does not occur.  Any finite
trajectory of height greater than one would contain a height-two tail, so no
other finite height occurs.  Finally, the number of parameters $p$ producing
$m\leq x$ is $O(\log x)$; division by the number of odd integers at most $x$
gives density zero for the height-one family.

The source states the classification and calls the height-two impossibility
easy, but omits the modular certificate
\cite[p.~1291]{Crandall1978}.  The calculation is reproduced by
\path{certificates/crandall_1978_checks.py}.
\end{proof}

\begin{nonclaim}
\label{nonclaim:Crandall-qxr-unboundedness}
Infinite height need not mean an unbounded orbit; it may arise from a
nontrivial cycle.  Crandall explicitly states that no unbounded trajectory
was then known for any $qx+1$ problem.  The $7x+1$ trajectory from $3$
exceeding $10^{2000}$ is a finite empirical excursion only.  The heuristic
drift $\log(q/4)>0$ for $q>3$ is not an unboundedness theorem
\cite[pp.~1291--1292]{Crandall1978}.
\end{nonclaim}
